\figura[0.38]{fig1.png}
{\centering \textit{Aquaspirillum magnetotacticum}\par}

\begin{apartat}{1}
El bacteri \textit{Aquaspirillum magnetotacticum} conté partícules molt petites, els magnetosomes, que són sensibles als camps magnètics. Fan servir el camp magnètic terrestre per a orientar-se en els oceans i nedar cap al pol Nord geogràfic. S'ha quantificat que una intensitat de camp magnètic inferior al $5\,\%$ del camp magnètic terrestre no té efectes sobre aquests bacteris. El camp magnètic terrestre és de $5,00\times 10^{-5}\ \mathrm{T}$. Si circula un corrent elèctric de $100\ \mathrm{A}$ per una línia submarina, a partir de quina distància d'aquesta línia el camp magnètic deixarà de tenir efecte sobre els bacteris? Considereu la línia submarina com un fil infinit i ignoreu els efectes de l'aigua del mar.
\end{apartat}

\begin{apartat}{1}
En la figura es mostren dos fils conductors rectilinis i infinitament llargs, que es troben situats als punts 1 i 2. Estan separats per $10,0\ \mathrm{m}$, són perpendiculars al pla del paper i per tots dos hi circula una mateixa intensitat de corrent de $100\ \mathrm{A}$ en el sentit que va cap endins del paper.

\figura[0.28]{fig2.png}

Representeu en un esquema el camp magnètic a la posició 1 generat pel conductor que passa per 2. Representeu també la força sobre el conductor que passa per 1 causada pel conductor que passa per 2, i calculeu el mòdul de la força que suporten $2,00\ \mathrm{m}$ del conductor que passa per 1.
\end{apartat}

\nota{El mòdul del camp magnètic a una distància $r$ d'un fil infinit pel qual circula una intensitat $I$ és $B = \dfrac{\mu_0 I}{2\pi r}$, en què $\mu_0 = 4\pi\times 10^{-7}\ \mathrm{T\,m\,A^{-1}}$.}
